\documentclass[10pt,a4paper]{amsart}
\usepackage{amsmath,amssymb,amsthm,graphicx}
\usepackage{verbatim}
\newenvironment{reference}{\comment}{\endcomment}
%\newenvironment{reference}{}{}
\newenvironment{slogan}{\comment}{\endcomment}
\newenvironment{history}{\comment}{\endcomment}

% For commutative diagrams we use Xy-pic
\usepackage[all]{xy}

% We use 2cell for 2-commutative diagrams.
\xyoption{2cell}
% Theorem environments.
%
\theoremstyle{plain}
\newtheorem{theorem}[subsection]{Theorem}
\newtheorem{proposition}[subsection]{Proposition}
\newtheorem{lemma}[subsection]{Lemma}

\theoremstyle{definition}
\newtheorem{definition}[subsection]{Definition}
\newtheorem{example}[subsection]{Example}
\newtheorem{exercise}[subsection]{Exercise}
\newtheorem{situation}[subsection]{Situation}

\theoremstyle{remark}
\newtheorem{remark}[subsection]{Remark}
\newtheorem{remarks}[subsection]{Remarks}

\numberwithin{equation}{subsection}

% Macros
%
\def\lim{\mathop{\mathrm{lim}}\nolimits}
\def\colim{\mathop{\mathrm{colim}}\nolimits}
\def\Spec{\mathop{\mathrm{Spec}}}
\def\Hom{\mathop{\mathrm{Hom}}\nolimits}
\def\Ext{\mathop{\mathrm{Ext}}\nolimits}
\def\SheafHom{\mathop{\mathcal{H}\!\mathit{om}}\nolimits}
\def\SheafExt{\mathop{\mathcal{E}\!\mathit{xt}}\nolimits}
\def\Sch{\mathit{Sch}}
\def\Mor{\mathop{\mathrm{Mor}}\nolimits}
\def\Ob{\mathop{\mathrm{Ob}}\nolimits}
\def\Sh{\mathop{\mathit{Sh}}\nolimits}
\def\NL{\mathop{N\!L}\nolimits}
\def\CH{\mathop{\mathrm{CH}}\nolimits}
\def\proetale{{pro\text{-}\acute{e}tale}}
\def\etale{{\acute{e}tale}}
\def\QCoh{\mathit{QCoh}}
\def\Ker{\mathop{\mathrm{Ker}}}
\def\Im{\mathop{\mathrm{Im}}}
\def\Coker{\mathop{\mathrm{Coker}}}
\def\Coim{\mathop{\mathrm{Coim}}}

% Boxtimes
%
\DeclareMathSymbol{\boxtimes}{\mathbin}{AMSa}{"02}

%
% Macros for moduli stacks/spaces
%
\def\QCohstack{\mathcal{QC}\!\mathit{oh}}
\def\Cohstack{\mathcal{C}\!\mathit{oh}}
\def\Spacesstack{\mathcal{S}\!\mathit{paces}}
\def\Quotfunctor{\mathrm{Quot}}
\def\Hilbfunctor{\mathrm{Hilb}}
\def\Curvesstack{\mathcal{C}\!\mathit{urves}}
\def\Polarizedstack{\mathcal{P}\!\mathit{olarized}}
\def\Complexesstack{\mathcal{C}\!\mathit{omplexes}}
% \Pic is the operator that assigns to X its picard group, usage \Pic(X)
% \Picardstack_{X/B} denotes the Picard stack of X over B
% \Picardfunctor_{X/B} denotes the Picard functor of X over B
\def\Pic{\mathop{\mathrm{Pic}}\nolimits}
\def\Picardstack{\mathcal{P}\!\mathit{ic}}
\def\Picardfunctor{\mathrm{Pic}}
\def\Deformationcategory{\mathcal{D}\!\mathit{ef}}
\usepackage{cleveref}
\crefname{originallemma}{original lemma}{original lemmas}
\Crefname{originallemma}{Original lemma}{Original lemmas}
\makeatletter
\newlabel{divisors-lemma-S1-no-embedded}{{0346}{source}{}{sourcepointer-1}{}}
\newlabel{divisors-lemma-S1-no-embedded@cref}{{[originallemma][0][]0346}{[1][0][]source}{}{}{}}
\makeatother
\crefname{originallemma}{original lemma}{original lemmas}
\Crefname{originallemma}{Original lemma}{Original lemmas}
\makeatletter
\newlabel{divisors-lemma-conormal-effective-Cartier-divisor}{{0B3P}{source}{}{sourcepointer-2}{}}
\newlabel{divisors-lemma-conormal-effective-Cartier-divisor@cref}{{[originallemma][0][]0B3P}{[1][0][]source}{}{}{}}
\makeatother
\crefname{originallemma}{original lemma}{original lemmas}
\Crefname{originallemma}{Original lemma}{Original lemmas}
\makeatletter
\newlabel{divisors-lemma-invertible-sheaf-sum-effective-Cartier-divisors}{{02OP}{source}{}{sourcepointer-3}{}}
\newlabel{divisors-lemma-invertible-sheaf-sum-effective-Cartier-divisors@cref}{{[originallemma][0][]02OP}{[1][0][]source}{}{}{}}
\makeatother
\crefname{originallemma}{original lemma}{original lemmas}
\Crefname{originallemma}{Original lemma}{Original lemmas}
\makeatletter
\newlabel{divisors-lemma-normal-effective-Cartier-divisor-S1}{{0B3S}{source}{}{sourcepointer-4}{}}
\newlabel{divisors-lemma-normal-effective-Cartier-divisor-S1@cref}{{[originallemma][0][]0B3S}{[1][0][]source}{}{}{}}
\makeatother
\crefname{originallemma}{original lemma}{original lemmas}
\Crefname{originallemma}{Original lemma}{Original lemmas}
\makeatletter
\newlabel{divisors-lemma-restriction-injective-open-contains-weakly-ass}{{0B3M}{source}{}{sourcepointer-5}{}}
\newlabel{divisors-lemma-restriction-injective-open-contains-weakly-ass@cref}{{[originallemma][0][]0B3M}{[1][0][]source}{}{}{}}
\makeatother
\crefname{originalreference}{original reference}{original references}
\Crefname{originalreference}{Original reference}{Original references}
\makeatletter
\newlabel{lemma-nice-meromorphic-function}{{0AX9}{source}{}{sourcepointer-6}{}}
\newlabel{lemma-nice-meromorphic-function@cref}{{[originalreference][0][]0AX9}{[1][0][]source}{}{}{}}
\makeatother
\crefname{originalreference}{original reference}{original references}
\Crefname{originalreference}{Original reference}{Original references}
\makeatletter
\newlabel{situation-vanishing}{{07UY}{source}{}{sourcepointer-7}{}}
\newlabel{situation-vanishing@cref}{{[originalreference][0][]07UY}{[1][0][]source}{}{}{}}
\makeatother
\crefname{originallemma}{original lemma}{original lemmas}
\Crefname{originallemma}{Original lemma}{Original lemmas}
\makeatletter
\newlabel{varieties-lemma-check-invertible-sheaf-trivial}{{0B40}{source}{}{sourcepointer-8}{}}
\newlabel{varieties-lemma-check-invertible-sheaf-trivial@cref}{{[originallemma][0][]0B40}{[1][0][]source}{}{}{}}
\makeatother
\crefname{originallemma}{original lemma}{original lemmas}
\Crefname{originallemma}{Original lemma}{Original lemmas}
\makeatletter
\newlabel{varieties-lemma-degree-tensor-product}{{0AYX}{source}{}{sourcepointer-9}{}}
\newlabel{varieties-lemma-degree-tensor-product@cref}{{[originallemma][0][]0AYX}{[1][0][]source}{}{}{}}
\makeatother
\title{The conormal sheaf of a nonempty Cartier divisor contained in the closed fibre is not trivial}
\author{The Stacks Project Authors}
\date{stacks-project:resolve.tex:local-snapshot-20261012; GNU FDL 1.2 or later}
\begin{document}
\maketitle
\noindent\textit{Editorial note.} All delivered mathematical text is the author's own, copied from the pinned Stacks Project chapter file resolve.tex; no statement or proof step was rewritten, reordered or removed.\par\medskip
\noindent\textit{Editorial note.} The counted proof is the author's complete argument: the choice of the component where the ratio of valuations is maximal and the permissible replacement of the divisor by a multiple, the use of the (S1) property to conclude that the chosen function vanishes on the divisor, and the degree computation showing that the resulting nonzero section of the conormal sheaf has a zero.\par\medskip
\noindent\textit{Editorial note.} The label lemma-nontrivial-normal-bundle names the counted statement, and the proof environment immediately following it in the source is delivered verbatim as its proof.\par\medskip
\noindent\textit{Editorial note.} The counted statement is carried by the author's own lemma environment in the chapter source.\par\medskip
\noindent\textit{Editorial note.} Pre-open gate: python operations/agent-authoring/tools/precard\_gate.py --source resolve.tex --label lemma-nontrivial-normal-bundle --start 1643 --end 1689 --local-id 28083 returned RESULT=OK (L28083 free, occupancy FREE, public interval clear). Pre-assembly tag self-check: canonical tag resolved and not present in reference/public-items.jsonl.\par\medskip
\noindent\textit{Editorial note.} The statement and proof are copied from the pinned chapter file; the cited inputs are kept as external pointers with their canonical Stacks tags.\par\medskip
\noindent\textit{Allowed standard prerequisites (editorial).}\begin{itemize}
\item Existence of a nice meromorphic function for the vanishing situation. Same chapter, the lemma producing a function with prescribed vanishing and pole orders along the components of the fibre. Supplies the function f used to measure the divisor and to produce the section whose zero forces nontriviality.
\item Relations between the conormal sheaf, the associated invertible sheaf and degrees. Divisors, the lemmas relating the invertible sheaf of a divisor to the conormal sheaf and to sums of divisors, and Varieties, the lemma computing the degree of a tensor product of invertible sheaves on a curve. Supplies the reinterpretations of the divisor and of its conormal sheaf used in the normalization step and in the degree computation.
\item The (S1) property of an effective Cartier divisor and absence of embedded associated points. Divisors, the lemma that a normal effective Cartier divisor satisfies (S1), and the lemmas that an (S1) scheme has no embedded associated points and that a section vanishing at the generic points of the components is zero. Supplies the conclusion that the image of the chosen function in the structure sheaf of the divisor is zero.
\item Nonvanishing of a section with a zero implies positive degree. Varieties, the lemma checking triviality of an invertible sheaf through degrees of sections on curves. Supplies the final conclusion that the conormal sheaf restricted to some component has positive degree.
\end{itemize}\medskip
\ifdefined\hypertarget\hypertarget{sourcepointer-1}{}\fi
\noindent\textit{Source reference.} Original lemma 0346: sources/stacks-chow/tags.txt, source lines 4081--4081. Prerequisite pointer invoked at proof line(s) 1681; it names a standard fact of a sibling chapter that the counted argument uses, and every step of the counted proof built on it is retained verbatim in the delivered body.\par\medskip
\ifdefined\hypertarget\hypertarget{sourcepointer-2}{}\fi
\noindent\textit{Source reference.} Original lemma 0B3P: sources/stacks-chow/tags.txt, source lines 13972--13972. Prerequisite pointer invoked at proof line(s) 1669; it names a standard fact of a sibling chapter that the counted argument uses, and every step of the counted proof built on it is retained verbatim in the delivered body.\par\medskip
\ifdefined\hypertarget\hypertarget{sourcepointer-3}{}\fi
\noindent\textit{Source reference.} Original lemma 02OP: sources/stacks-chow/tags.txt, source lines 3524--3524. Prerequisite pointer invoked at proof line(s) 1668; it names a standard fact of a sibling chapter that the counted argument uses, and every step of the counted proof built on it is retained verbatim in the delivered body.\par\medskip
\ifdefined\hypertarget\hypertarget{sourcepointer-4}{}\fi
\noindent\textit{Source reference.} Original lemma 0B3S: sources/stacks-chow/tags.txt, source lines 13975--13975. Prerequisite pointer invoked at proof line(s) 1678; it names a standard fact of a sibling chapter that the counted argument uses, and every step of the counted proof built on it is retained verbatim in the delivered body.\par\medskip
\ifdefined\hypertarget\hypertarget{sourcepointer-5}{}\fi
\noindent\textit{Source reference.} Original lemma 0B3M: sources/stacks-chow/tags.txt, source lines 13970--13970. Prerequisite pointer invoked at proof line(s) 1682; it names a standard fact of a sibling chapter that the counted argument uses, and every step of the counted proof built on it is retained verbatim in the delivered body.\par\medskip
\ifdefined\hypertarget\hypertarget{sourcepointer-6}{}\fi
\noindent\textit{Source reference.} Original reference 0AX9: sources/stacks-resolve/pinned-original-source.tex, source lines 1623--1623. Prerequisite pointer invoked at proof line(s) 1656; it names a standard fact of the same chapter that the counted argument uses, and every step of the counted proof built on it is retained verbatim in the delivered body.\par\medskip
\ifdefined\hypertarget\hypertarget{sourcepointer-7}{}\fi
\noindent\textit{Source reference.} Original reference 07UY: sources/stacks-resolve/pinned-original-source.tex, source lines 1584--1584. Prerequisite pointer invoked at proof line(s) 1645; it names a standard fact of the same chapter that the counted argument uses, and every step of the counted proof built on it is retained verbatim in the delivered body.\par\medskip
\ifdefined\hypertarget\hypertarget{sourcepointer-8}{}\fi
\noindent\textit{Source reference.} Original lemma 0B40: sources/stacks-chow/tags.txt, source lines 13983--13983. Prerequisite pointer invoked at proof line(s) 1688; it names a standard fact of a sibling chapter that the counted argument uses, and every step of the counted proof built on it is retained verbatim in the delivered body.\par\medskip
\ifdefined\hypertarget\hypertarget{sourcepointer-9}{}\fi
\noindent\textit{Source reference.} Original lemma 0AYX: sources/stacks-chow/tags.txt, source lines 13805--13805. Prerequisite pointer invoked at proof line(s) 1671; it names a standard fact of a sibling chapter that the counted argument uses, and every step of the counted proof built on it is retained verbatim in the delivered body.\par\medskip
\begin{lemma}
\label{lemma-nontrivial-normal-bundle}
In Situation \ref{situation-vanishing} assume $X$ is normal.
Let $Z \subset X$ be a nonempty effective Cartier divisor such that
$Z \subset X_s$ set theoretically.
Then the conormal sheaf of $Z$ is not trivial.
More precisely, there exists an $i$ such that $C_i \subset Z$
and $\deg(\mathcal{C}_{Z/X}|_{C_i}) > 0$.
\end{lemma}

\begin{proof}
We will use the observations made following Situation \ref{situation-vanishing}
without further mention. Let $f$ be a function as in
Lemma \ref{lemma-nice-meromorphic-function}.
Let $\xi_i \in C_i$ be the generic point. Let
$\mathcal{O}_i$ be the local ring of $X$ at $\xi_i$. Then $\mathcal{O}_i$
is a discrete valuation ring. Let $e_i$ be the valuation of
$f$ in $\mathcal{O}_i$, so $e_i > 0$. Let $h_i \in \mathcal{O}_i$ be a local
equation for $Z$ and let $d_i$ be its valuation. Then $d_i \geq 0$.
Choose and fix $i$ with $d_i/e_i$ maximal (then $d_i > 0$ as
$Z$ is not empty). Replace $f$ by $f^{d_i}$ and $Z$ by $e_iZ$.
This is permissible, by the relation
$\mathcal{O}_X(e_i Z) = \mathcal{O}_X(Z)^{\otimes e_i}$,
the relation between the conormal sheaf and $\mathcal{O}_X(Z)$
(see Divisors, Lemmas
\ref{divisors-lemma-invertible-sheaf-sum-effective-Cartier-divisors}
and \ref{divisors-lemma-conormal-effective-Cartier-divisor}, and
since the degree gets multiplied by $e_i$, see
Varieties, Lemma \ref{varieties-lemma-degree-tensor-product}.
Let $\mathcal{I}$ be the ideal sheaf of $Z$ so that
$\mathcal{C}_{Z/X} = \mathcal{I}|_Z$. Consider the image $\overline{f}$
of $f$ in $\Gamma(Z, \mathcal{O}_Z)$. By our choices above we see
that $\overline{f}$ vanishes in the generic points of irreducible
components of $Z$ (these are all generic points of $C_j$ as $Z$ is
contained in the special fibre). On the other hand, $Z$ is $(S_1)$ by
Divisors, Lemma \ref{divisors-lemma-normal-effective-Cartier-divisor-S1}.
Thus the scheme $Z$ has no embedded associated points and
we conclude that $\overline{f} = 0$ (Divisors, Lemmas
\ref{divisors-lemma-S1-no-embedded} and
\ref{divisors-lemma-restriction-injective-open-contains-weakly-ass}).
Hence $f$ is a global section of $\mathcal{I}$
which generates $\mathcal{I}_{\xi_i}$ by construction.
Thus the image $s_i$ of $f$ in $\Gamma(C_i, \mathcal{I}|_{C_i})$ is nonzero.
However, our choice of $f$ guarantees that $s_i$ has a zero at $x_i$.
Hence the degree of $\mathcal{I}|_{C_i}$ is $> 0$ by
Varieties, Lemma \ref{varieties-lemma-check-invertible-sheaf-trivial}.
\end{proof}

\end{document}
